\chapter[The Right Hand]{How to Use the Right Hand / Playing the Open Strings}

In this chapter, we will explore the various capabilities of the right hand, starting with proper posture and mechanics.

\section{The job of the right hand}

The right hand is responsible for several different essential aspects of guitar playing. Perhaps most obviously, the right hand is in charge of pulling or striking the strings in order for the guitar to actually be heard. Most simply, the right hand can pull one string at a time, using either the thumb, index, middle, or ring finger. Later, you will learn more complex skills like using multiple fingers at the same time, or using the right hand to strum chord patterns. The right hand is also responsible for creating different rhythms in music, whether by plucking or strumming. The vast majority of rhythmic concepts will be practiced using the right hand fingers.

\section{Right hand ``posture" and mechanics}

Our first step is to ensure that the entire right arm is placed correctly on the guitar. Remember from Chapter 3 that the right forearm contacts the top right ``corner" of the guitar, directly above the saddle. Once placed, the left hand is adjusted so that it rests above the sound hole. To start, practice pushing the palm down against the strings, then raise the fingers and make a little ``tent" where the palm is off of the strings.

Once you get comfortable with this, reset the arm and try setting up the tent without placing the palm on the strings first. Then, start putting your right fingers on specific strings:

\begin{itemize}
	
	\item Thumb goes on the top metal string (closest to your face)
	
	\item Index finger goes on the top clear plastic string (the thickest plastic string)
	
	\item Middle finger goes on the middle clear plastic string
	
	\item Ring fingers goes on the bottom clear plastic string (the thinnest plastic string)
	
\end{itemize}

\begin{center}
	\fbox{\textbf{Note: }Make sure your right hand wrist is not touching the wood of the guitar.}
\end{center}

% Image of hand tent with fingers on correct strings

\section{Push towards the wood}

As you prepare to play your first notes, avoid a common beginner mistake of pulling the string away from the guitar. Instead, think of each string like a key or button that must be depressed. Push the string down towards the sound hole, then release and let the string ring. 

Set up your hand tent with the aforementioned string placement and prepare to use each finger to play the string it's touching. 

\begin{itemize}
	
	\item Thumb pushes through the top metal string, landing on the next metal string below
	
	\item Index finger pushes through the top clear nylon string, landing on the next metal string above
	
	\item Middle finger pushes through the middle clear nylon string, landing on the top nylon string (where index finger starts)
	
	\item Ring finger pushes through the bottom clear nylon string, landing on the middle nylon string (where middle finger starts) 
	
\end{itemize} 

Repeat this process, slowly, several times over. Make sure you are always landing on a different string directly after playing a string with any finger.

\section{Right hand finger letters}

From this point on, we need to learn the shorthand system for referring to the right hand fingers, known as \textit{pima}. This will save us a lot of time in the future when we can refer to each finger with a single letter: \textit{p}, \textit{i}, \textit{m}, or \textit{a}. The letters come from the Spanish words for each of the four fingers used. They will always appear in music as lower case and italicized.

\begin{itemize}
	
	\item \textit{p}, from \textit{pulgar} (thumb)
	
	\item \textit{i}, from \textit{indice} (index finger)
	
	\item \textit{m}, from \textit{medio dedo} (middle finger)
	
	\item \textit{a}, from \textit{anular} (ring finger)
	
\end{itemize}

% Image of the right hand with pima labeled

\begin{figure}[h]
	\begin{center}
		\large
		\lilypondfile{basic_pima_example.ly}
		\caption{\textit{A simple example of how right hand letters look in sheet music.}}
	\end{center}
\end{figure}

\section{The great nail debate}

\section{Open strings of the guitar}

As we start to acclimate ourselves to playing the strings with the right hand fingers, we should also go ahead and learn the letters and numbers used to refer to each of the six strings. 

\begin{itemize}
	\item The bottom nylon string is the 1st string or the E string
	\item The middle nylon string is the 2nd string or the B string
	\item The top nylon string is the 3rd string or the G string
	\item The bottom metal string is the 4th string or the D string
	\item The middle metal string is the 5th string or the A string
	\item The top metal string is the 6th string or the E string
\end{itemize}

Notice how the 1st and 6th strings share the same letter, E. Because of this, you must specify which E string you are talking about. You can refer to the 1st string as the nylon E string, and the 6th string as the metal E string. 

You will play the open strings so many times during class that you will most likely memorize their letters and numbers quickly, however there is a famous mnemonic trick used to memorize the letters even faster. \textbf{Warning!} If you use this trick, remember that the letters are written in reverse order (6-5-4-3-2-1).

\begin{center}
	\fbox{Eddie Ate Dynamite, Good Bye Eddie}
\end{center}

At this point, you may have noticed the term \textit{open strings} used earlier. \textit{Open strings} refers to playing any of the six strings with the right hand fingers without using the left hand fingers to stop the strings on a fret. We will play our open strings almost every single day, because it allows us to focus on creating a beautiful sound without adding the complexity of using the left hand at the same time.

\section{The first open string warmup}

The open string warmup is a basic set of actions meant to train you to use proper technique in creating a beautiful sound on the guitar. 

\noindent \textbf{Before You Begin}
\begin{itemize}
	\item Check and make sure your seated posture is correct, going over each of the four main contact points.
	\item Rest your left hand on your left leg, or on the bottom left corner of the guitar body.
	\item Set up your right hand ``tent" over the strings, this time placing the \textit{i} finger on the 1st string E and the \textit{p} finger on the 4th string D.
\end{itemize}

\noindent \textbf{Open String Warmup Part 1}
\begin{itemize}
	\item Set a metronome to 4/4 time and 60 BPM.
	\item Play the 1st string E with the \textit{i} finger four times, trying to play at the same time as the beat.
	\item Repeat the above step on the 2nd string B and the 3rd string G.
	\item Play the 4th string D with the \textit{p} finger four times, again trying to match the speed of the beat.
	\item Repeat the above step on the 5th string A and the 6th string E.
	\item After having played each string four times, rest for four beats.
	\item Play each string four times in reverse: E-A-D-G-B-E. Make sure to use the \textit{p} finger for the 6th, 5th, and 4th strings, and the \textit{i} finger for the 3rd, 2nd, and 1st strings.
\end{itemize}

At this point, you will have played each string four times from 1 to 6, and each string again four times in reverse from 6 to 1. This is considered the first round or cycle of the warmup. The next two rounds follow a very similar pattern, but with different rhythms played in the left hand.

\newpage

\noindent \textbf{Open String Warmup Part 2}
\begin{itemize}
	\item Play each open string two times, letting the string ring for two beats each time. Play the strings in order from 1st to 6th, rest for four beats, then play in reverse order from 6th to 1st.
	\item Play each open string one time, letting the string ring for four beats each time. Play the strings in the same order as above, resting four beats before the change in direction.
\end{itemize}

Although you can play the entire open string warmup by following written or verbal directions, we can go ahead and take a look at what a sheet music version of this warmup would look like.

\begin{figure}
	\begin{center}
		\lilypondfile{open_string_warmup_basic.ly}
		\caption{\textit{The basic open string warmup written out in sheet music notation.}}
	\end{center}
\end{figure}